\documentclass[12pt]{article}

\usepackage{times}
\usepackage{graphicx}
\usepackage{subfigure}
\usepackage{color}
\usepackage{url}
\usepackage{latexsym}
\usepackage{geometry}
\usepackage{indentfirst}
\usepackage{array}
\usepackage{tabularx}
\usepackage{xtab}
\usepackage{fancyvrb}
\usepackage[center, bf]{caption}
\usepackage{listings}
\usepackage{url}
\usepackage{placeins}

\begin{document}
  \begin{titlepage}
    \centering
     \begin{figure}[ht]
       \centering
       \includegraphics[angle=-90, keepaspectratio=true, width=8cm]{logo.eps}
     \end{figure}
     {
     \large \bfseries Nomadic Comunication\par
     AA 2008-2009
     }\par
     \vspace{1.5cm}
     {
     \Large \bfseries \textcolor{blue}{Report 1: The tile!} \par
     }
     \vspace{1.0cm}
     {
     \large \bfseries {Group N. ###} \par
     }
     \vspace{0.3cm}
     {\large \bfseries {Author 1, author 2, etc...}
     \vspace{1.0cm}
\begin{abstract}
	%Abstract here.
\end{abstract}
\end{titlepage}

\thispagestyle{empty}
\tableofcontents
\clearpage
\setcounter{page}{1}

%%%%
% Usefull things in latex:
% itemize, enumerate, lstlistings (to report source code, shell output, whatever), emph, tables
%
%%%
% TODO
% - Refer to Professor template for some usefull hints: http://disi.unitn.it/locigno/index.php/teaching-duties/nomadic-communications/aa08-09
% - *Always* use \cite and \ref when needed. Adding it when the text is already written is senseless waste of time, the result will be hugly.
% - Use \footnote if necessary, it can make things more readable.
% - Please indent your Latex code! Use tabs, not spaces.
% - English! Write in grammatically correct English. (Aspell can check latex source files)
% - This is a formal document that requires formal writing. This means that "it's", "I've", "ain't" and others must be replaced by their extended form.
% - Use \emph the first time you call something new or when you introduce a new abbreviation i.e.: Internet Protocol (\emph{IP}). Do not abbreviate access point to AP in descriptions and similar
%

\section{Introduction} \label{intro}
	% Describe what we are gonna do, the general ideas and what we are gonna talk about.
	% Say what is needed to understand the sections that will follow.

\section{Theoretical background} \label{theory}
	% Everything that comes from theory and that is needed to understand and describe our work.
	% A simple description of how an IEEE 802.11 protocol works to introduce the fact that there is some overhead added by the protocol.

	\subsection{Differences between 802.11b and 802.11g} \label{theory:prot_differences}
		% Yes, I know it's unbelivable but there are some differences, for real! :)
		% Calculate protocols upperbound limitations and similar things (report the formulas!).
		% DO NOT USE PROFESSOR SLIDES AS SOURCE OF INFORMATIONS SINCE MOST OF VALUES ARE COMPLETELY WRONG!!!!!!!!!
		% Download the IEEE standard from here: http://disi.unitn.it/locigno/didattica/NC/802.11-2007.pdf and make a large use of ^F.

\section{Testbed Setup} \label{setup}
	% Say what is the testbed used for tests including *everything* necessary, both hardware, software, hardware position, air humidity and everything meaningful for the tests.
	%
	% We have used 2 pc (client C and server S), 1 access point (AP) and 1 switch. The idea is simple:
	% - AP, C and S are connected to the switch on a wired LAN
	% - we use 2 separate class C LANs (wlan and clan):
	%      one for wired traffic like retriving tcpdump and iperf outputs
	%      the other one for the wireless tests
	% - C connects to S via SSH over clan, starts iperf server and starts tcpdump
	% - C uses iperf to send UDP traffic to S over AP on wlan
	% - C gets out outputs from S using clan
	\subsection{Hardware} \label{setup:hardware}
		% mainly C wifi card configuration, ap settings

	\subsection{Software} \label{setup:software}
		% ssh keys, iwconfig, python software, iperf, tcpdump, gnuplot, whatever used

\subsection{Parameters} \label{parameters}
	% Parameters of the tests. Everything that influences the tests and the results we are trying to look at. This includes both static and variables things.
	% Wifi channel speed used, fragmentation threshold, rst threshold and similars

\section{Results} \label{results}
	% Our results and something that underlines that we understand the mean of the results. Results are not self evident!
	% Any formula used to calculate the efficency or whatever must be written and described.
	\subsection{IEEE 802.11b} \label{resutls:11b}
	\subsection{IEEE 802.11g} \label{resutls:11g}

\section{Conclusion} \label{concluion}
	% Conclusion must say what we understand from our results, what is the learned lesson and so on.

\addcontentsline{toc}{section}{Bibliography}
\bibliographystyle{plain}
\bibliography{biblio} % put bibliografy in a file called biblio.bib
% Include all books, manuals, softwares home pages etc used for the tests.
% Including links to wikipedia is like saying that we are a bunch of fools... better not to :)
\nocite{*}
\appendix % if necessary
% We should do something for Andrea's software...
\end{document}
